\textbf{Whom We will Involve.} Phase~I (complete) involved five participants assigned female sex at birth, three with confirmed endometriosis. Sessions used semi-structured think-aloud usability methods with the interactive prototype and the synthetic patient dataset described in Appendix~D. Phase~II (proposed) will enroll $N=160$ adult Phendo users (80 per arm) with confirmed endometriosis and a scheduled gynecology visit, recruited through the existing Phendo and Citizen Endo recruitment infrastructure. Phase~III is outlined but not the current commitment.

\textbf{Risks and Benefits.} Reflecting on chronic illness can bring up difficult emotions: pain, frustration, grief, and memories of dismissal by providers. Sharing health information also carries privacy risks, especially for a condition that is stigmatized and historically dismissed. Anticipated benefits include the creation of a space in which the patient's account is treated as authoritative; support for the cognitive and emotional preparation that patients already undertake before clinical visits; and, if Phase~II is successful, evidence that conversational-speech-based representations are a feasible alternative to structured-form self-tracking for an enigmatic disease.

\textbf{Participant Protections.} Participants will provide informed consent and may withdraw at any time without penalty. EndoChron stores all health data on-device; patients may delete it at any time by clearing local storage or uninstalling the app. Aggregate engagement metrics are shared only upon consent. Audio is discarded after transcription. Data submitted to upstream LLMs for summarization will be free of personally identifying information: requests will \textit{only entail transcripts} and not include name, age, gender, location, device signatures, or any other personally identifiable information. For Phase~II, an adverse event surveillance item will be administered at the post-visit timepoint asking whether the app caused distress, intrusive worry, or other negative affect. Phendo engagement data accessed for stratification will be governed by the existing IRB-approved data-use agreement. Phendo stratification data are accessed under the existing data-use agreement and de-identified.

\textbf{IRB.} All procedures will be reviewed by the Columbia IRB under an amendment to the existing Phendo protocol (\href{https://pmc.ncbi.nlm.nih.gov/articles/PMC7314826/}{\texttt{AAAQ9812}}), consistent with prior Phendo research.

\textbf{Inclusion.} Endometriosis primarily affects people with female reproductive anatomy. We will recruit women and welcome non-binary, transgender, and intersex people with the condition. Phase~II will exclude participants under 18, consistent with the existing Phendo protocol's pediatric provisions. Race, ethnicity, education, and socioeconomic status will be measured at baseline and reported transparently. The existing Phendo cohort is known to over-represent white, college-educated, English-speaking participants\cite{Urteaga2022Engagement}. Phase~III is planned in part to address that limitation.

\textbf{Broader Ethics.} Endometriosis patients are routinely dismissed by the healthcare system and historically under-studied. Following the citizen-science approach of prior Phendo work\cite{McKillop2016Participatory}, we will treat participants as collaborators rather than subjects, center their interpretations as the gold standard against which algorithmic representations are checked, and share findings back with the patient community. EndoChron is intentionally designed to complement clinical care and human relationships, not to replace them. No chatbot or simulated-empathy features are included. No attempt is made to override participants' narratives at any temporal resolution. We position the summarization AI as a reflective anchor and explicitly state this in the UI with explicit calls-to-action that all AI-generated content is editable by the user at any time.


\newpage
